% ===== PRÉAMBULE CANONIQUE — à coller VERBATIM en tête de CHAQUE .tex =====
\documentclass[11pt,a4paper]{article}
\usepackage[utf8]{inputenc}\usepackage[T1]{fontenc}\usepackage{lmodern}
\usepackage[french]{babel}
\usepackage{amsmath,amssymb,mathtools}\usepackage{icomma}
\usepackage{siunitx}\sisetup{locale=FR}  % \qty{}{}, \si{}, \ang{}
\usepackage{xcolor}\usepackage{colortbl}
\usepackage{tikz}\usepackage{pgfplots}\pgfplotsset{compat=1.18}
\usepackage[siunitx,europeanresistors]{circuitikz} % circuits (élec.)
\usepackage[most]{tcolorbox}\usepackage{enumitem}\usepackage{array,booktabs}
\usepackage[a4paper,margin=2.2cm]{geometry}
\usetikzlibrary{arrows.meta,calc,angles,quotes,positioning,patterns,%
  backgrounds,shapes.geometric,decorations.pathmorphing,decorations.markings,intersections}
\definecolor{primary}{RGB}{222,49,99}\definecolor{primarydark}{RGB}{150,22,55}
\definecolor{defcol}{RGB}{0,114,178}\definecolor{methcol}{RGB}{52,92,148}
\definecolor{excol}{RGB}{0,135,90}\definecolor{attncol}{RGB}{200,40,55}
\tcbset{boxbase/.style={enhanced,breakable,boxrule=0.9pt,arc=2.5pt,
  left=3mm,right=3mm,top=2.3mm,bottom=2.3mm,fonttitle=\bfseries,coltitle=white}}
% --- boîtes TUTORIEL (noms EXACTS, ne pas renommer) ---
\newtcolorbox{cle}[1][]{boxbase,colback=defcol!6,colframe=defcol,
  title={\textsc{À retenir}\ifx\relax#1\relax\relax\else\textmd{ : \itshape #1}\fi}}
\newtcolorbox{methode}[1][]{boxbase,colback=methcol!7,colframe=methcol,sharp corners,
  title={\textsc{Méthode}\ifx\relax#1\relax\relax\else\textmd{ --- \itshape #1}\fi}}
\newtcolorbox{exemplebox}[1][]{boxbase,colback=excol!5,colframe=excol,
  title={\textsc{Exemple}\ifx\relax#1\relax\relax\else\textmd{ : \itshape #1}\fi}}
\newtcolorbox{piege}[1][]{boxbase,colback=attncol!6,colframe=attncol,
  title={\textsc{Piège}\ifx\relax#1\relax\relax\else\textmd{ : \itshape #1}\fi}}
\newtcolorbox{remarque}[1][]{boxbase,colback=black!4,colframe=black!45,
  title={\textsc{Remarque}\ifx\relax#1\relax\relax\else\textmd{ : \itshape #1}\fi}}
% --- boîtes EXERCICE / CORRIGÉ (noms EXACTS, auto-comptées) ---
\newtcolorbox[auto counter]{exo}[1][]{boxbase,colback=primary!4,colframe=primary,
  title={\textsc{Exercice}~\thetcbcounter\ifx\relax#1\relax\relax\else\textmd{ --- \itshape #1}\fi}}
\newtcolorbox[auto counter]{corrige}[1][]{boxbase,colback=excol!4,colframe=excol,
  title={\textsc{Corrigé}~\thetcbcounter\ifx\relax#1\relax\relax\else\textmd{ --- \itshape #1}\fi}}
\newcommand{\vv}[1]{\vec{#1}}\newcommand{\ds}{\displaystyle}
\newcommand{\R}{\mathbb{R}}\newcommand{\N}{\mathbb{N}}\newcommand{\Z}{\mathbb{Z}}
% (un \usepackage{fancyhdr}+en-tête est possible ; il est ignoré côté web)
% =========================================================================
\begin{document}
\sisetup{per-mode=symbol}

\begin{corrige}[vrai ou faux sur un graphe de vitesse]
\begin{enumerate}[label=\alph*)]
  \item \textbf{V} : la vitesse croît linéairement (droite de pente $\frac{6}{2}=\qty{3}{\metre\per\second\squared}$), c'est un MRUA.
  \item \textbf{F} : la vitesse est constante \emph{non nulle} ($\qty{6}{\metre\per\second}$) : le mobile avance en MRU, il n'est pas à l'arrêt.
  \item \textbf{V} : à $t=\qty{6}{\second}$ la vitesse passe de positive à négative en s'annulant : le mobile fait demi-tour.
  \item \textbf{V} : pente $=\dfrac{-3-6}{7-4}=\dfrac{-9}{3}=\qty{-3}{\metre\per\second\squared}$.
  \item \textbf{V} : aire du triangle $\tfrac12\times2\times6=\qty{6}{\metre}$.
\end{enumerate}
\textbf{Bilan : 4 affirmations vraies} (a, c, d, e) ; seule b est fausse.
\end{corrige}

\begin{corrige}[décrypter une équation horaire parabolique]
\begin{enumerate}[label=\alph*)]
  \item $v(t)=-2t+6$ (\si{\metre\per\second}) et $a(t)=\qty{-2}{\metre\per\second\squared}$ : accélération constante $\Rightarrow$ mouvement rectiligne uniformément varié (ici décéléré tant que $v>0$).
  \item $v=0 \Leftrightarrow t=\qty{3}{\second}$ ; la trajectoire atteint alors son point le plus avancé ($x_{\max}=x(3)=-9+18-5=\qty{4}{\metre}$), c'est le sommet de la parabole.
  \item $x=0 \Leftrightarrow t^2-6t+5=0 \Leftrightarrow (t-1)(t-5)=0$ : $t=\qty{1}{\second}$ et $t=\qty{5}{\second}$.
  \item Avant $t=\qty{3}{\second}$, $v>0$ : il avance dans le sens de $Ox$ ; après, $v<0$ : il recule.
\end{enumerate}
\end{corrige}

\begin{corrige}[déplacement contre distance parcourue]
\begin{enumerate}[label=\alph*)]
  \item $\ds \Delta x=\int_0^3(-2t+4)\,\mathrm{d}t=\big[-t^2+4t\big]_0^3=-9+12=\qty{3}{\metre}$.
  \item $v=0 \Leftrightarrow -2t+4=0 \Leftrightarrow t=\qty{2}{\second}$ (avant : $v>0$ ; après : $v<0$).
  \item On sépare les deux phases (distance $=\int|v|$) :
        \[ \int_0^2(-2t+4)\,\mathrm{d}t+\int_2^3(2t-4)\,\mathrm{d}t
           =\big[-t^2+4t\big]_0^2+\big[t^2-4t\big]_2^3=4+1=\qty{5}{\metre}. \]
        La distance ($\qty{5}{\metre}$) dépasse le déplacement ($\qty{3}{\metre}$) car le mobile est revenu en arrière.
\end{enumerate}
\end{corrige}

\begin{corrige}[affirmations de cours à trancher]
\begin{enumerate}[label=\alph*)]
  \item \textbf{V} : $a=\text{cte}\Rightarrow v(t)=v_0+at$, fonction affine du temps.
  \item \textbf{F} : au sommet d'un tir vertical $v=0$ mais $a=-g\neq0$ ; vitesse nulle n'implique pas accélération nulle.
  \item \textbf{V} : c'est la définition de l'intégrale $\int v\,\mathrm{d}t$.
  \item \textbf{F} : si $v<0$ (mobile reculant), une accélération négative le fait \emph{accélérer} dans le sens négatif. Il ne ralentit que si $a$ et $v$ sont de signes opposés.
  \item \textbf{V} : le déplacement est nul, donc $v_{\text{moy}}=\Delta x/\Delta t=0$ (la distance, elle, n'est pas nulle).
\end{enumerate}
\textbf{Bilan : 3 affirmations vraies} (a, c, e).
\end{corrige}

\end{document}
